\paragraph{\texorpdfstring{$S_5\supset S_4\supset A_4\supset V_4$}{S5>S4>A4>V4} 子群层的分歧剖分与中间曲线的显式模型}
沿用定理 \ref{thm:xi-terminal-zm-delta-map-hurwitz-passport-s5} 的次数 \(5\) 覆盖
\(\pi_\delta:E\to\PP^1_\delta\) 与结论 \ref{con:xi-terminal-zm-delta-s5-regular-closure-genus-109} 的正则闭包
\[
L/\QQ(\delta),\qquad \Gal(L/\QQ(\delta))\cong S_5,\qquad \QQ(X)=L,
\]
并在自然五点作用下取一点稳定子 \(H\simeq S_4\le S_5\)，从而识别
\[
\QQ(E)=L^{H},\qquad E\simeq X/H,
\]
并记相应商映射为 \(\pi:X\to E\)。

\begin{conclusion}[\texorpdfstring{$S_4$}{S4}-层分歧剖分与分歧点数 \(18\)]\label{con:xi-terminal-zm-delta-s4-layer-branch-decomposition-18}
设 \(\infty\in\PP^1_\delta\) 为 \(\delta\) 的无穷远分歧值，有限分歧值集合为
\[
\{\delta=4\}\ \cup\ \{\delta\in\overline{\QQ}:\ B_5(\delta)=0\}
\]
（定理 \ref{thm:xi-terminal-zm-delta-map-hurwitz-passport-s5}）。
对任意有限分歧值 \(t\)，其纤维在 \(E\) 上分解为
\[
\delta^{-1}(t)=2P_t+Q_{t,1}+Q_{t,2}+Q_{t,3}
\]
（按零点重数计），其中 \(P_t\) 为 \(\pi_\delta\) 的简单分歧点（分歧指数 \(2\)），其余三点 \(Q_{t,i}\) 在 \(\pi_\delta\) 下不分歧。
令
\[
B:=\sum_{t}\ \sum_{i=1}^{3} Q_{t,i},\qquad \deg(B)=18.
\]
则：
\begin{enumerate}
  \item \(\pi:X\to E\) 为 \(S_4\)-伽罗瓦覆叠，\(\deg(\pi)=|S_4|=24\)；
  \item \(\pi\) 在无穷远纤维的唯一点 \(O\in E\) 上方不分歧；
  \item \(\pi\) 的分歧点集恰为 \(B\) 的支撑：对每个 \(Q\in \mathrm{Supp}(B)\)，其惯性群为 \(S_4\) 中由换位生成的阶 \(2\) 循环群，且
  \[
  \#\pi^{-1}(Q)=12,\qquad e_P(\pi)=2\quad(\forall P\in\pi^{-1}(Q));
  \]
  \item \(\deg(R_\pi)=216\)，并与 \(g(X)=109\)、\(g(E)=1\) 的 Hurwitz 公式严格相符。
\end{enumerate}
\end{conclusion}
\begin{proof}
设 \(v\) 为 \(\QQ(\delta)\) 上的分歧赋值，对应惯性群阶 \(|I_v|\in\{5,2\}\)。
取 \(w\) 为 \(\QQ(E)\) 上位于 \(v\) 之上的赋值，则塔式分歧指数满足
\[
e(L/\QQ(E);w)=\frac{e(L/\QQ(\delta);v)}{e(\QQ(E)/\QQ(\delta);w)}
=\frac{|I_v|}{e(E/\PP^1_\delta;w)}.
\]
当 \(v\) 对应无穷远分歧点时，\(|I_v|=5\) 且 \(\pi_\delta\) 在 \(O\) 处全分歧 \(e(E/\PP^1_\delta;O)=5\)，故 \(e(L/\QQ(E);O)=1\)，从而 \(\pi\) 在 \(O\) 处不分歧。
\par
当 \(v\) 对应有限分歧值 \(t\) 时，\(|I_v|=2\)。
在纤维 \(\delta^{-1}(t)\) 中仅有 \(P_t\) 满足 \(e(E/\PP^1_\delta;P_t)=2\)，故 \(e(L/\QQ(E);P_t)=1\)；
而对其余三点 \(Q_{t,i}\) 有 \(e(E/\PP^1_\delta;Q_{t,i})=1\)，从而 \(e(L/\QQ(E);Q_{t,i})=2\)。
因此 \(\pi\) 的分歧点恰为全部 \(Q_{t,i}\)，总数为 \(6\cdot 3=18\)。
由于 \(\pi\) 为 \(S_4\)-伽罗瓦覆叠且每个分歧点惯性阶为 \(2\)，故分歧点之上点数为 \(24/2=12\)，且每点贡献 \(e-1=1\)，从而 \(\deg(R_\pi)=18\cdot 12=216\)。
代入 Hurwitz 公式
\(
2g(X)-2=24(2g(E)-2)+\deg(R_\pi)
\)
即得到 \(g(X)=109\)，与结论 \ref{con:xi-terminal-zm-delta-s5-regular-closure-genus-109} 一致。证毕。
\end{proof}

\begin{conclusion}[主除子恒等式与 \texorpdfstring{$B$}{B} 的解析切出]\label{con:xi-terminal-zm-delta-branch-complement-principal-divisor}
令
\[
R_{\delta,\mathrm{fin}}:=\sum_{t} P_t
\]
为六个有限分歧点之和（按几何点计一次），\(O\in E\) 为 \(\delta\) 的极点。
则：
\begin{enumerate}
  \item 分歧除子满足线性等价 \(R_{\delta,\mathrm{fin}}\sim 6O\)，从而在椭圆曲线群律下 \(\sum_t P_t=O\)；
  \item 微分的除子满足严格恒等式
  \[
  \operatorname{div}(\dd\delta)=R_{\delta,\mathrm{fin}}-6O;
  \]
  \item 令 \(G:=(\delta-4)B_5(\delta)\in\QQ(E)\)，并定义
  \[
  \Theta:=\frac{G}{(\dd\delta)^2}\in\QQ(E),
  \]
  则 \(\Theta\) 的除子为
  \[
  \operatorname{div}(\Theta)=B-18O.
  \]
  特别地，\(\Theta\) 在 \(B\) 上均为一次零点，且在 \(O\) 处有唯一极点阶 \(18\)；并且 \(B\sim 18O\)、\(\sum_{Q\in \mathrm{Supp}(B)}Q=O\)。
\end{enumerate}
\end{conclusion}
\begin{proof}
由 \(K_E\sim 0\) 与 \(K_{\PP^1}\sim -2(\infty)\)，对 \(\pi_\delta:E\to\PP^1_\delta\) 的 Hurwitz 公式给出
\[
R_\delta\sim 2\cdot \pi_\delta^*(\infty)=10O.
\]
又 \(O\) 处全分歧指数为 \(5\)，故
\(
R_\delta=4O+R_{\delta,\mathrm{fin}}
\)，从而 \(R_{\delta,\mathrm{fin}}\sim 6O\)，并在 \(\mathrm{Pic}^0(E)\cong E\) 下等价于 \(\sum_t P_t=O\)。
\par
取 \(\PP^1\) 的局部坐标 \(t\)，则 \(\operatorname{div}(\dd t)=-2(\infty)\)。
由微分拉回公式
\[
\operatorname{div}(\dd\delta)=\pi_\delta^*\operatorname{div}(\dd t)+R_\delta=-2\pi_\delta^*(\infty)+R_\delta=-10O+(4O+R_{\delta,\mathrm{fin}})=R_{\delta,\mathrm{fin}}-6O.
\]
\par
对每个有限分歧值 \(t\)，有
\(
\operatorname{div}(\delta-t)=2P_t+Q_{t,1}+Q_{t,2}+Q_{t,3}-5O
\)。
对 \(t=4\) 与 \(B_5\) 的五个根相乘得到
\[
\operatorname{div}(G)=2R_{\delta,\mathrm{fin}}+B-30O.
\]
于是
\[
\operatorname{div}(\Theta)=\operatorname{div}(G)-2\operatorname{div}(\dd\delta)
=(2R_{\delta,\mathrm{fin}}+B-30O)-2(R_{\delta,\mathrm{fin}}-6O)=B-18O.
\]
由 \(\operatorname{div}(\Theta)\) 的度为零得 \(B\sim 18O\)，并在群律下得到 \(\sum_{Q\in \mathrm{Supp}(B)}Q=O\)。证毕。
\end{proof}

\begin{conclusion}[\texorpdfstring{$A_4$}{A4} 交叉子群与规范化纤维积：属 \(10\) 的二次模型]\label{con:xi-terminal-zm-delta-a4-fiberproduct-y10-explicit}
令 \(A_5\trianglelefteq S_5\) 为交错群，并记
\[
C_{A_5}:=X/A_5,\qquad C_{A_5}:\ s^2=(\delta-4)B_5(\delta)
\]
（结论 \ref{con:xi-terminal-zm-delta-a5-fixedfield-hyperelliptic-genus2}）。
取一点稳定子 \(H\simeq S_4\) 并令 \(A_4:=A_5\cap H\)。
置
\[
Y_{10}:=X/A_4.
\]
则
\[
Y_{10}\ \simeq\ \mathrm{Norm}\bigl(E\times_{\PP^1_\delta} C_{A_5}\bigr),
\]
并且在 \(E\) 上具有显式二次模型
\[
Y_{10}:\ z^2=(\delta-4)B_5(\delta).
\]
此外，二次覆叠 \(Y_{10}\to E\) 的分歧点集恰为 \(B\)，从而 \(g(Y_{10})=10\)。
\end{conclusion}
\begin{proof}
在 \(S_5\)-伽罗瓦扩张 \(L/\QQ(\delta)\) 中，对任意子群 \(U,V\le S_5\) 有域论恒等式
\(
L^{U\cap V}=L^{U}\cdot L^{V}
\)（合成取于 \(L\) 内）。
取 \(U=A_5\)、\(V=H\)，并注意 \(A_5\cap H=A_4\)，得
\[
L^{A_4}=L^{A_5}\cdot L^{H}=\QQ(C_{A_5})\cdot \QQ(E),
\]
其对应的曲线模型即为纤维积 \(E\times_{\PP^1_\delta} C_{A_5}\) 的规范化。
将 \(C_{A_5}\) 的方程在纤维积中取 \(\delta\) 的拉回，即得到 \(z^2=(\delta-4)B_5(\delta)\) 的二次模型。
\par
由结论 \ref{con:xi-terminal-zm-delta-branch-complement-principal-divisor}，
\(
\operatorname{div}(G)=2R_{\delta,\mathrm{fin}}+B-30O
\)，故 \(G\) 在 \(B\) 上为奇阶零点而在 \(R_{\delta,\mathrm{fin}}\) 上为偶阶零点；
因此二次扩张 \(\QQ(E)(\sqrt{G})/\QQ(E)\) 的分歧点集恰为 \(B\)。
由 Hurwitz 公式
\[
2g(Y_{10})-2=2(2g(E)-2)+\deg(B)=18
\]
得到 \(g(Y_{10})=10\)。证毕。
\end{proof}

\begin{conclusion}[\texorpdfstring{$V_4$}{V4}-resolvent：\texorpdfstring{$S_3$}{S3} 伽罗瓦覆叠与三次无分歧循环上覆]\label{con:xi-terminal-zm-delta-v4-s3-resolvent-z28-etale-c3}
令 \(V_4\trianglelefteq H\simeq S_4\) 为 Klein 四元子群，并置
\[
Z_{28}:=X/V_4.
\]
则 \(Z_{28}\to E\) 为 \(S_3\)-伽罗瓦覆叠（由短正合列 \(1\to V_4\to S_4\to S_3\to 1\) 给出），次数为 \(6\)，且其分歧点集恰为 \(B\)，每个分歧点的惯性群在 \(S_3\) 中为换位生成的阶 \(2\) 子群。
设 \(A_3\trianglelefteq S_3\) 为唯一指数 \(2\) 正规子群，则其逆像为 \(A_4\)，从而 \(Z_{28}\) 的唯一二次中间商为 \(Y_{10}=X/A_4\)，并且自然映射
\[
Z_{28}\longrightarrow Y_{10}
\]
为次数 \(3\) 的无分歧循环覆叠（伽罗瓦群 \(C_3\)）。
\end{conclusion}
\begin{proof}
由 \(V_4\trianglelefteq H\) 且 \(H/V_4\simeq S_3\) 可知 \(Z_{28}\to E\) 为 \(S_3\)-伽罗瓦覆叠且次数为 \(6\)。
在结论 \ref{con:xi-terminal-zm-delta-s4-layer-branch-decomposition-18} 中，\(\pi:X\to E\) 的惯性在 \(H\) 内为换位；
其在商群 \(H/V_4\simeq S_3\) 中仍为换位，故 \(Z_{28}\to E\) 的分歧点集与 \(\pi\) 相同，均为 \(B\)。
\par
在 \(S_3\) 中唯一指数 \(2\) 正规子群为 \(A_3\)；
其在 \(H\) 中的逆像为 \(A_4\)，从而 \(Z_{28}/A_3\simeq X/A_4=Y_{10}\)。
由于 \(A_3\) 为 \(S_3\) 的换位换位积生成的 \(3\)-循环子群，其不含阶 \(2\) 元素，故在 \(B\) 上惯性在该层必为平凡，从而 \(Z_{28}\to Y_{10}\) 无分歧且为循环三次覆叠。证毕。
\end{proof}

\begin{conclusion}[三条共轭三次子覆叠与 \texorpdfstring{$S_3$}{S3} 闭包]\label{con:xi-terminal-zm-delta-three-cubic-subcovers-s3-closure}
在 \(S_3\) 中换位子群共有三条共轭类。
取 \(D_8\le H\simeq S_4\) 为 \(V_4\) 的一个上覆，使得 \(D_8/V_4\) 对应于 \(S_3\) 的某一换位子群（从而 \(|D_8|=8\) 且 \([H:D_8]=3\)）。
置
\[
C_{10}:=X/D_8.
\]
则：
\begin{enumerate}
  \item \(C_{10}\to E\) 为次数 \(3\) 的非伽罗瓦覆叠，在 \(B\) 上处处为简单分歧（局部分歧指数 \(2\)），且无其他分歧点；因此 \(g(C_{10})=10\)；
  \item 该覆叠的伽罗瓦闭包为 \(Z_{28}\to E\)；等价地，\(Z_{28}\) 为 \(C_{10}\to E\) 的 \(S_3\)-伽罗瓦闭包；
  \item 在 \(E\) 上存在恰好三条此类次数 \(3\) 覆叠，它们在 \(Z_{28}\) 的 \(S_3\)-作用下相互置换。
\end{enumerate}
\end{conclusion}
\begin{proof}
由 \([H:D_8]=3\) 得 \(C_{10}\to E\) 的次数为 \(3\)。
在 \(B\) 上，\(H\) 的惯性为换位，其在陪集集 \(H/D_8\) 上诱导的置换型为 \(2+1\)，从而 \(C_{10}\to E\) 在 \(B\) 上简单分歧且无额外分歧点。
Hurwitz 公式给出
\[
2g(C_{10})-2=3(2g(E)-2)+\deg(B)=18,
\]
故 \(g(C_{10})=10\)。
\par
由于 \(D_8/V_4\) 为 \(S_3\) 的换位子群，其共轭闭包生成 \(S_3\)；
在 \(H\) 中相应的共轭闭包生成 \(V_4\) 之上全体 \(H\)，故函数域意义下的伽罗瓦闭包为 \(X/V_4=Z_{28}\)。
换位子群在 \(S_3\) 中恰有三条共轭类，故得到三条互为共轭的次数 \(3\) 子覆叠，并被 \(S_3\) 置换。证毕。
\end{proof}

\begin{conclusion}[标准四维等型分量由五个椭圆商生成且与 \texorpdfstring{$E^4$}{E^4} 同源]\label{con:xi-terminal-zm-delta-jacobian-standard-part-e4}
令 \(H_i\le S_5\)（\(i=1,\dots,5\)）为自然五点作用下的五个点稳定子（两两共轭且同构于 \(S_4\)），并记
\[
E_i:=X/H_i,\qquad \pi_i:X\to E_i.
\]
由 \(S_5\) 在 \(X\) 上的甲板变换，曲线 \(\{E_i\}\) 在 \(\QQ\) 上两两同构，且同构于既定椭圆曲线 \(E\)。
令
\[
\iota_i:\ E_i\simeq J(E_i)\xrightarrow{\ \pi_i^*\ }J(X)
\]
为诱导同态，并令
\[
A_{\mathrm{std}}:=\langle\,\iota_1(E_1),\dots,\iota_5(E_5)\,\rangle\subset J(X)
\]
为其生成的最小阿贝尔子簇。
则 \(A_{\mathrm{std}}\) 的维数为 \(4\)，并存在 \(\QQ\)-同源
\[
E^5/\Delta(E)\ \sim_{\QQ}\ A_{\mathrm{std}},
\qquad
 A_{\mathrm{std}}\ \sim_{\QQ}\ E^4,
\]
其中 \(\Delta(E)\subset E^5\) 为对角嵌入。
并且在结论 \ref{con:xi-terminal-zm-delta-s5-jacobian-qisogeny-decomposition} 的记号下，有
\(
A_{\mathrm{std}}\sim_{\QQ}A_{(4,1)}
\)。
\end{conclusion}
\begin{proof}
取每个 \(E_i\) 上非零正则微分 \(\omega_i\)，则 \(\pi_i^*\omega_i\in H^0(X,\Omega_X^1)\) 非零。
由于 \(H_i\) 在 \(S_5\) 中两两共轭，且 \(S_5\) 作用置换这些商映射，故 \(\{\pi_i^*\omega_i\}_{i=1}^{5}\) 的线性张成空间在 \(S_5\) 下同构于五点置换表示。
另一方面，\(H^0(X,\Omega_X^1)^{S_5}\cong H^0(X/S_5,\Omega^1)=H^0(\PP^1_\delta,\Omega^1)=0\)，故该置换表示中平凡不变向量不出现，于是其非平凡部分为标准表示 \(V_{(4,1)}\)，维数为 \(4\)。
因此由 \(\iota_i\) 生成的阿贝尔子簇 \(A_{\mathrm{std}}\) 的正则微分空间对偶于该 \(4\) 维子表示，从而 \(\dim A_{\mathrm{std}}=4\)，并与 \(A_{(4,1)}\) 等型一致。
\par
定义同态
\[
\Phi:\ E_1\times\cdots\times E_5\longrightarrow J(X),\qquad
(P_1,\dots,P_5)\longmapsto \sum_{i=1}^5 \iota_i(P_i).
\]
其像为 \(A_{\mathrm{std}}\)。
对角子簇 \(\Delta(E)\) 的像在微分层面对应置换表示的平凡方向；由于该方向在 \(H^0(X,\Omega^1)\) 中不出现，故 \(\Delta(E)\subset\ker\Phi\)，从而 \(\Phi\) 下降为同态
\[
\overline{\Phi}:\ E^5/\Delta(E)\longrightarrow A_{\mathrm{std}}.
\]
两端维数同为 \(4\)，且 \(\overline{\Phi}\) 在微分上为同构，故其为同源；再由 \(E^5/\Delta(E)\sim_{\QQ}E^4\) 得结论。证毕。
\end{proof}

\endinput
